% Все настройки шаблона находятся здесь.
% Файл рассчитан на XeLaTeX или LuaLaTeX; pdfLaTeX не поддерживается.

% Метаданные документа — меняйте их для нового конспекта.
\newcommand{\coursename}{Математический анализ}
\newcommand{\notestitle}{Конспект лекций}
\newcommand{\studentname}{Автор: Николай Сухов}
\newcommand{\yearname}{2026}
\newcommand{\sourcenote}{%
  Формулировки: В.~А.~Зорич, «Математический анализ», часть~I,
  10-е изд., 2019. © В.~А.~Зорич и правообладатели.\\
  Конспект предназначен для личного использования.}

% Русский и английский языки.
\usepackage{fontspec}
\usepackage{polyglossia}
\setdefaultlanguage{russian}
\setotherlanguage{english}

% Эти шрифты входят в Windows и содержат кириллицу. Поэтому шаблон не зависит
% от того, зарегистрировал ли MiKTeX дополнительные шрифты Libertinus.
\setmainfont{Cambria}
\setsansfont{Segoe UI}
\setmonofont{Consolas}[Scale=MatchLowercase]
\newfontfamily\cyrillicfont{Cambria}
\newfontfamily\cyrillicfontsf{Segoe UI}
\newfontfamily\cyrillicfonttt{Consolas}[Scale=MatchLowercase]

% Математика. mathtools должен загружаться раньше unicode-math.
\usepackage{mathtools}
\usepackage{amsthm}
\usepackage{unicode-math}
\setmathfont{Cambria Math}

\newcommand{\N}{\mathbb{N}}
\newcommand{\Z}{\mathbb{Z}}
\newcommand{\Q}{\mathbb{Q}}
\newcommand{\R}{\mathbb{R}}
% \C уже определена в LaTeX как команда для акцента, поэтому используем \CC.
\newcommand{\CC}{\mathbb{C}}
\newcommand{\eps}{\varepsilon}
\newcommand{\dd}{\mathop{}\!\mathrm{d}}
\newcommand{\diff}[2]{\frac{\dd #1}{\dd #2}}
\newcommand{\pdiff}[2]{\frac{\partial #1}{\partial #2}}
\newcommand{\abs}[1]{\left\lvert #1\right\rvert}
\newcommand{\norm}[1]{\left\lVert #1\right\rVert}
\newcommand{\inner}[2]{\left\langle #1,#2\right\rangle}
\newcommand{\set}[1]{\left\{#1\right\}}
\DeclareMathOperator{\sgn}{sgn}
\DeclareMathOperator{\grad}{grad}
\DeclareMathOperator{\rank}{rank}
\DeclareMathOperator{\card}{card}
\newcommand{\powerset}[1]{\mathcal{P}\!\left(#1\right)}

% Страница и типографика.
\usepackage[a4paper,margin=25mm,headheight=21pt]{geometry}
\usepackage{microtype}
\usepackage{setspace}
\setstretch{1.08}
\setlength{\parindent}{1.25em}
\setlength{\parskip}{0.25em}
\usepackage{xcolor}
\definecolor{accent}{HTML}{1F4E79}
\usepackage{enumitem}
\setlist{nosep,leftmargin=2em}
\usepackage{booktabs}
\usepackage{array}
\usepackage{graphicx}
\graphicspath{{images/}}

% Теоремы, определения и другие смысловые блоки.
% Пример: \begin{theorem}{Название}{label} ... \end{theorem}
\usepackage[most]{tcolorbox}
\tcbset{
  framedstatement/.style={
    enhanced,
    before={\par},
    after={\par},
    colback=white,
    colframe=accent,
    colbacktitle=white,
    coltitle=accent,
    coltext=black,
    boxrule=0.6pt,
    arc=0mm,
    left=3mm,right=3mm,top=2mm,bottom=2mm,
    before skip=8pt,after skip=8pt,
    fonttitle=\bfseries
  },
  plainstatement/.style={
    enhanced,
    breakable,
    blanker,
    before={\par},
    after={\par},
    left=0mm,right=0mm,top=1mm,bottom=1mm,
    before skip=7pt,after skip=7pt,
    fonttitle=\bfseries,
    coltitle=accent,
    coltext=black
  }
}

\newtcbtheorem[number within=section]{theorem}{Теорема}{
  framedstatement
}{thm}
\newtcbtheorem[use counter from=theorem]{lemma}{Лемма}{
  framedstatement
}{lem}
\newtcbtheorem[use counter from=theorem]{proposition}{Утверждение}{
  plainstatement
}{prop}
\newtcbtheorem[use counter from=theorem]{axiom}{Аксиома}{
  plainstatement
}{ax}
\newtcbtheorem[use counter from=theorem]{definition}{Определение}{
  plainstatement
}{def}
\newtcbtheorem[use counter from=theorem]{corollary}{Следствие}{
  plainstatement
}{cor}
\newtcbtheorem[use counter from=theorem]{properties}{Свойства}{
  plainstatement
}{prop-list}
\newtcbtheorem[use counter from=theorem]{example}{Пример}{
  plainstatement
}{ex}

\newtcolorbox{remark}{
  plainstatement,title=Замечание
}
\newtcolorbox{important}{
  plainstatement,title=Важно
}

\renewcommand{\proofname}{\textcolor{accent}{Доказательство}}

% Ссылки, заголовки и колонтитулы.
\usepackage{hyperref}
\hypersetup{
  colorlinks=true,
  linkcolor=accent,
  citecolor=accent,
  urlcolor=accent,
  pdftitle={\coursename{} — \notestitle},
  pdfauthor={\studentname}
}
\usepackage[capitalise,noabbrev]{cleveref}
\crefname{equation}{формула}{формулы}
\Crefname{equation}{Формула}{Формулы}

\usepackage{fancyhdr}
\pagestyle{fancy}
\fancyhf{}
\renewcommand{\chaptermark}[1]{%
  \markboth{Глава~\thechapter.\enspace #1}{}%
}
\renewcommand{\sectionmark}[1]{%
  \markright{\S\,\thesection\enspace #1}%
}
\fancyhead[L]{\scriptsize\color{accent}Глава~\thechapter}
\fancyhead[R]{\scriptsize\color{accent}\nouppercase{\rightmark}}
\fancyfoot[C]{\thepage}
\renewcommand{\headrulewidth}{0.4pt}

% На первой странице главы оставляем только номер страницы.
\fancypagestyle{plain}{%
  \fancyhf{}
  \fancyfoot[C]{\thepage}
  \renewcommand{\headrulewidth}{0pt}
}

\usepackage{titlesec}
\titleformat{\chapter}[display]
  {\LARGE\bfseries\color{accent}\raggedright}{Глава~\thechapter}{0.5em}{}
\titlespacing*{\chapter}{0pt}{-20pt}{30pt}
\titleformat{\section}
  {\large\bfseries\color{accent}}{\S\,\thesection}{0.75em}{}
\titleformat{\subsection}
  {\large\bfseries\color{accent}}{\thesubsection}{0.75em}{}

% Параграф без даты в колонтитуле и на полях. Необязательный первый аргумент
% задаёт сокращённое название только для колонтитула, полное остаётся в тексте
% и оглавлении: \lecture[Короткое название]{Полное название}.
\newcommand{\lecture}[2][]{%
  \if\relax\detokenize{#1}\relax
    \section{#2}%
  \else
    \begingroup
      \renewcommand{\sectionmark}[1]{%
        \markright{\S\,\thesection\enspace #1}%
      }%
      \section{#2}%
    \endgroup
  \fi
}
